\section{COMPLEXES OF A-MODULES}

In this paragraph, A denotes a ring. When we speak of A-modules without further specification, we shall always mean left A-modules.

By graded modules we shall mean graded modules of type Z (II, p.~164). If M is a graded A-module, with grading \((\mathbf{M}_n)_{n \in \mathbf{Z}}\) , we set \(\mathbf{M}^n = \mathbf{M}_{-n}\) and we say that \(\mathbf{M}_n\) (resp. \(\mathbf{M}^n\) ) is the homogeneous component of descending degree \(n\) (resp. of ascending degree \(n\) ) of M. If \(u: \mathbf{M} \to \mathbf{N}\) is a graded homomorphism of degree \(p\) of graded A-modules (II, p.~166), we denote by \(u_n: \mathbf{M}_n \to \mathbf{M}_{n+p}\) (resp. \(u^n: \mathbf{M}^n \to \mathbf{M}^{n-p}\) ) the homomorphism deduced from \(u\) ; it is called the homogeneous component of descending degree (resp. ascending degree) \(n\) of \(u\) ; we also say that \(u\) is of descending degree \(p\) or of ascending degree \(-p\) .

